% GENERATED FILE. Edit canonical lessons, then run npm run generate:books.
% canonical-source-hash: fnv1a64:e77a89e12e2bd564
% canonical-lessons: GU-C377-dasa, GU-C377-sthir, GU-C377-samanya, GU-C377-badlavu, GU-C377-vanasvu

\chapter{As Things Stand}
\label{ch:gu-as-things-stand}
\hlchaptermodality{377}

\begin{chapteropening}
\textbf{By the end of this chapter:} \emph{I can say what we will do if the situation changes.}

Complete the last lesson of chapter 377: I can say what we will do if the situation changes.
\end{chapteropening}

\section[\texorpdfstring{\gu{દશા} (daśā)}{દશા (daśā)}]{\gu{દશા} (daśā) — a plight; a state of affairs}

\label{lesson:GU-C377-dasa}



\begin{quote}
Before the new one: say the Gujarati for slackness, then the Gujarati for to get better.
\end{quote}

\subsection*{You'll want to know: \gu{દશા}}
\textbf{\gu{દશા}} — \emph{daśā} — \textquotedblleft{}a plight; a state of affairs\textquotedblright{}.

\begin{quote}
\gu{જો} \gu{આ} \gu{વરસે} \gu{પણ} \gu{વરસાદ} \gu{ઓછો} \gu{પડશે} \gu{તો} \gu{ખેતરોની} \gu{દશા} \gu{ખરાબ} \gu{થશે}.
\end{quote}

\begin{quote}
\emph{jo ā varse paṇ varsād ocho paḍśe to khetronī daśā kharāb thaśe.}
\end{quote}

\begin{quote}
If the rain falls short again this year, the fields will be in a bad way.
\end{quote}

A feminine noun.

\begin{cousinweb}
\textbf{\gu{દશા}} (\emph{daśā}) is a tatsama, borrowed from Sanskrit \emph{daśā}, \textquotedblleft{}a state, a condition of life\textquotedblright{} (Monier-Williams).
\end{cousinweb}

\begin{grammarlens}[title={What you've built}]
One more word for saying what happens if.
\end{grammarlens}

\subsection*{Your turn}
\begin{itemize}
\raggedright
  \item \emph{Say it:} \emph{daśā}
  \item \emph{Say it:} \emph{daśā}, once more
  \item \emph{Say it:} \emph{jo ā varse paṇ varsād ocho paḍśe to khetronī daśā kharāb thaśe.} — If the rain falls short again this year, the fields will be in a bad way.
  \item \emph{Recall:} say the Gujarati for slackness, then the Gujarati for to get better, then say \emph{daśā} again
\end{itemize}

\begin{tcolorbox}[breakable,colback=teal!4,colframe=teal!35!black,title={Before you move on}]
What does \gu{દશા} mean? (\textquotedblleft{}A plight; a state of affairs\textquotedblright{}.) Say it once more.
\end{tcolorbox}
\section[\texorpdfstring{\gu{સ્થિર} (sthir)}{સ્થિર (sthir)}]{\gu{સ્થિર} (sthir) — steady}

\label{lesson:GU-C377-sthir}



\begin{quote}
Before the new one: say the Gujarati for to get better, then the Gujarati for a plight.
\end{quote}

\subsection*{You'll want to know: \gu{સ્થિર}}
\textbf{\gu{સ્થિર}} — \emph{sthir} — \textquotedblleft{}steady\textquotedblright{}.

\begin{quote}
\gu{જો} \gu{ભાવ} \gu{સ્થિર} \gu{રહેશે} \gu{તો} \gu{વેપારીઓ} \gu{ખુશ} \gu{થશે}.
\end{quote}

\begin{quote}
\emph{jo bhāv sthir raheśe to vepārīo khuś thaśe.}
\end{quote}

\begin{quote}
If prices stay steady, the traders will be happy.
\end{quote}

\begin{cousinweb}
\textbf{\gu{સ્થિર}} (\emph{sthir}) is a tatsama, borrowed from Sanskrit \emph{sthira}, \textquotedblleft{}firm, hard, steady\textquotedblright{} (Monier-Williams).
\end{cousinweb}

\begin{grammarlens}[title={What you've built}]
One more word for saying what happens if.
\end{grammarlens}

\subsection*{Your turn}
\begin{itemize}
\raggedright
  \item \emph{Say it:} \emph{sthir}
  \item \emph{Say it:} \emph{sthir}, once more
  \item \emph{Say it:} \emph{jo bhāv sthir raheśe to vepārīo khuś thaśe.} — If prices stay steady, the traders will be happy.
  \item \emph{Recall:} say the Gujarati for to get better, then the Gujarati for a plight, then say \emph{sthir} again
\end{itemize}

\begin{tcolorbox}[breakable,colback=teal!4,colframe=teal!35!black,title={Before you move on}]
What does \gu{સ્થિર} mean? (\textquotedblleft{}Steady\textquotedblright{}.) Say it once more.
\end{tcolorbox}
\section[\texorpdfstring{\gu{સામાન્ય} (sāmānya)}{સામાન્ય (sāmānya)}]{\gu{સામાન્ય} (sāmānya) — normal; ordinary}

\label{lesson:GU-C377-samanya}



\begin{quote}
Before the new one: say the Gujarati for a plight, then the Gujarati for steady.
\end{quote}

\subsection*{You'll want to know: \gu{સામાન્ય}}
\textbf{\gu{સામાન્ય}} — \emph{sāmānya} — \textquotedblleft{}normal; ordinary\textquotedblright{}.

\begin{quote}
\gu{જો} \gu{બધું} \gu{સામાન્ય} \gu{રહેશે} \gu{તો} \gu{અમે} \gu{સોમવારે} \gu{પાછા} \gu{આવીશું}.
\end{quote}

\begin{quote}
\emph{jo badhũ sāmānya raheśe to ame somvāre pāchā āvīśũ.}
\end{quote}

\begin{quote}
If all goes normally, we will be back on Monday.
\end{quote}

\begin{cousinweb}
\textbf{\gu{સામાન્ય}} (\emph{sāmānya}) is a tatsama, borrowed from Sanskrit \emph{sāmānya}, \textquotedblleft{}equal, alike, common\textquotedblright{} (Monier-Williams).
\end{cousinweb}

\begin{grammarlens}[title={What you've built}]
One more word for saying what happens if.
\end{grammarlens}

\subsection*{Your turn}
\begin{itemize}
\raggedright
  \item \emph{Say it:} \emph{sāmānya}
  \item \emph{Say it:} \emph{sāmānya}, once more
  \item \emph{Say it:} \emph{jo badhũ sāmānya raheśe to ame somvāre pāchā āvīśũ.} — If all goes normally, we will be back on Monday.
  \item \emph{Recall:} say the Gujarati for a plight, then the Gujarati for steady, then say \emph{sāmānya} again
\end{itemize}

\begin{tcolorbox}[breakable,colback=teal!4,colframe=teal!35!black,title={Before you move on}]
What does \gu{સામાન્ય} mean? (\textquotedblleft{}Normal; ordinary\textquotedblright{}.) Say it once more.
\end{tcolorbox}
\section[\texorpdfstring{\gu{બદલાવું} (badlāvũ)}{બદલાવું (badlāvũ)}]{\gu{બદલાવું} (badlāvũ) — to become different}

\label{lesson:GU-C377-badlavu}



\begin{quote}
Before the new one: say the Gujarati for steady, then the Gujarati for normal.
\end{quote}

\subsection*{You'll want to know: \gu{બદલાવું}}
\textbf{\gu{બદલાવું}} — \emph{badlāvũ} — \textquotedblleft{}to become different\textquotedblright{}.

\begin{quote}
\gu{જો} \gu{હવામાન} \gu{બદલાશે} \gu{તો} \gu{અમે} \gu{ઘરે} \gu{રહીશું}.
\end{quote}

\begin{quote}
\emph{jo havāmān badlāśe to ame ghare rahīśũ.}
\end{quote}

\begin{quote}
If the weather changes, we will stay at home.
\end{quote}

Past (it …): \gu{તે} \gu{શહેર} \gu{ઘણું} \gu{બદલાયું}, \emph{te śaher ghaṇũ badlāyũ}.

\begin{cousinweb}
\textbf{\gu{બદલાવું}} (\emph{badlāvũ}) is the \emph{-āvũ} partner of \textbf{\gu{બદલવું}} (\emph{badalvũ}), \textquotedblleft{}to change\textquotedblright{}, which is usually traced to Arabic \emph{badal}, \textquotedblleft{}a substitute, an exchange\textquotedblright{}, through Persian.
\end{cousinweb}

\begin{grammarlens}[title={What you've built}]
One more word for saying what happens if.
\end{grammarlens}

\subsection*{Your turn}
\begin{itemize}
\raggedright
  \item \emph{Say it:} \emph{badlāvũ}
  \item \emph{Say it:} \emph{badlāvũ}, once more
  \item \emph{Say it:} \emph{jo havāmān badlāśe to ame ghare rahīśũ.} — If the weather changes, we will stay at home.
  \item \emph{Recall:} say the Gujarati for steady, then the Gujarati for normal, then say \emph{badlāvũ} again
\end{itemize}

\begin{tcolorbox}[breakable,colback=teal!4,colframe=teal!35!black,title={Before you move on}]
What does \gu{બદલાવું} mean? (\textquotedblleft{}To become different\textquotedblright{}.) Say it once more.
\end{tcolorbox}
\section[\texorpdfstring{\gu{વણસવું} (vaṇasvũ)}{વણસવું (vaṇasvũ)}]{\gu{વણસવું} (vaṇasvũ) — to worsen}

\label{lesson:GU-C377-vanasvu}



\begin{quote}
Before the new one: say the Gujarati for normal, then the Gujarati for to become different.
\end{quote}

\subsection*{You'll want to know: \gu{વણસવું}}
\textbf{\gu{વણસવું}} — \emph{vaṇasvũ} — \textquotedblleft{}to worsen\textquotedblright{}.

\begin{quote}
\gu{જો} \gu{સારવાર} \gu{મોડી} \gu{થશે} \gu{તો} \gu{દર્દીની} \gu{તબિયત} \gu{વણસશે}.
\end{quote}

\begin{quote}
\emph{jo sārvār moḍī thaśe to dardīnī tabiyat vaṇasśe.}
\end{quote}

\begin{quote}
If treatment is delayed, the patient's health will worsen.
\end{quote}

Past (it …): \gu{તે} \gu{દર્દીની} \gu{તબિયત} \gu{વણસી}, \emph{te dardīnī tabiyat vaṇsī}.

\begin{cousinweb}
\textbf{\gu{વણસવું}} (\emph{vaṇasvũ}) is inherited from Sanskrit \emph{vínaśyati}, \textquotedblleft{}perishes\textquotedblright{}; CDIAL lists the Gujarati verb, \textquotedblleft{}to be spoilt\textquotedblright{} (CDIAL 11770).
\end{cousinweb}

\begin{grammarlens}[title={What you've built}]
One more word for saying what happens if.
\end{grammarlens}

\subsection*{Your turn}
\begin{itemize}
\raggedright
  \item \emph{Say it:} \emph{vaṇasvũ}
  \item \emph{Say it:} \emph{vaṇasvũ}, once more
  \item \emph{Say it:} \emph{jo sārvār moḍī thaśe to dardīnī tabiyat vaṇasśe.} — If treatment is delayed, the patient's health will worsen.
  \item \emph{Recall:} say the Gujarati for normal, then the Gujarati for to become different, then say \emph{vaṇasvũ} again
\end{itemize}

\begin{tcolorbox}[breakable,colback=teal!4,colframe=teal!35!black,title={Before you move on}]
What does \gu{વણસવું} mean? (\textquotedblleft{}To worsen\textquotedblright{}.) Say it once more.
\end{tcolorbox}
